\documentclass[conference]{IEEEtran}

\usepackage{iftex}
\ifPDFTeX
\usepackage[utf8]{inputenc}
\usepackage[T1]{fontenc}
\fi
\usepackage{amsmath,bm}
\usepackage{newtxtext,newtxmath}
\usepackage{textcomp}
\interdisplaylinepenalty=2500
\usepackage{array,booktabs}
\usepackage{mathtools}
\usepackage{graphicx}
\usepackage{pgf}
\usepackage{pgfplots}
\usepackage{siunitx}
\usepackage[final]{microtype}
\usepackage{cite}
\usepackage[breaklinks=true,hidelinks]{hyperref}
\usepackage{bookmark}
\usepackage[nameinlink,capitalise,noabbrev]{cleveref}

\DeclareRobustCommand{\vct}[1]{\bm{#1}}
\DeclareRobustCommand{\mat}[1]{\bm{#1}}
\DeclareRobustCommand{\matg}[1]{\bm{#1}}
\DeclareMathOperator{\clip}{clip}
\DeclareMathOperator{\wrap}{wrap}
\DeclareMathOperator{\sgn}{sgn}
\newcommand{\R}{\mathbb{R}}
\newcommand{\E}{\mathbb{E}}
\newcommand{\pct}[1]{#1\%}
\providecommand{\mathdefault}[1]{#1}
\pgfplotsset{compat=1.18}

\sisetup{detect-weight=true,detect-family=true,per-mode=symbol}
\AtBeginEnvironment{equation}{\small}
\AtBeginEnvironment{align}{\small}
\emergencystretch=2em
\allowdisplaybreaks

\pdfstringdefDisableCommands{%
    \def\SI#1#2{#1 #2}%
    \def\bm#1{#1}%
    \def\vct#1{#1}%
    \def\mat#1{#1}%
    \def\matg#1{#1}%
    \def\degree{deg}%
    \def\omega{omega}%
    \def\theta{theta}%
}

\title{Single-IMU Wave-Direction and Apparent Travel-Sense Estimation with Adaptive Frequency Tracking (draft)}
\author{%
    \IEEEauthorblockN{Mikhail S. Grushinskiy}
    \IEEEauthorblockA{Independent Researcher, USA, Fair Lawn\\
    Bareboat Necessities GitHub Project\\
    Email: mgrouch@users.github.com}%
}

\begin{document}

\IfFileExists{w3d-ou-validation-macros-generated.tex-part}{%
  \input{w3d-ou-validation-macros-generated.tex-part}%
}{}
\providecommand{\OUValidationDirectionAbsError}{--}
\providecommand{\OUValidationDirectionRMSE}{--}
\providecommand{\OUValidationDirectionWorstAbsError}{--}
\providecommand{\OUValidationDirectionDominant}{--}
\providecommand{\OUValidationDirectionUncertain}{--}
\providecommand{\OUValidationTravelCorrect}{--}
\providecommand{\OUValidationTravelWrong}{--}
\providecommand{\OUValidationTravelUnresolved}{--}
\providecommand{\OUValidationTravelAbsError}{--}

\maketitle

\begin{abstract}
A real-time estimator is presented for recovering the dominant horizontal wave
propagation axis and resolving its $180^\circ$ travel-sense ambiguity from a
single marine IMU.  A Kalman-updated adaptive notch supplies a slowly varying
carrier frequency.  Horizontal inertial acceleration is expressed in a levelled
heading frame and fitted with sine and cosine quadratures, making the recovered
axis insensitive to arbitrary carrier-phase offset.  A second stage compares
signed acceleration along that axis with the vertical orbital quadrature to
resolve apparent propagation sense.  In a ten-seed simulation study spanning
four directional JONSWAP and four PM--Stokes seas, the pooled absolute
circular-mean axis error is \OUValidationDirectionAbsError{} degrees and no
scenario exceeds \OUValidationDirectionWorstAbsError{} degrees.  The
truth-referenced travel-sense decision is correct for
\OUValidationTravelCorrect\% of samples, wrong for
\OUValidationTravelWrong\%, and unresolved for
\OUValidationTravelUnresolved\%.  The direct estimate is an encounter,
boat-relative direction; conversion to intrinsic or ground-referenced wave
direction requires vessel velocity, current, and a dispersion model.
\end{abstract}

\begin{IEEEkeywords}
wave direction, inertial measurement unit, IMU, marine sensing, adaptive notch
filter, Kalman filter, wave buoy, directional waves, travel sense
\end{IEEEkeywords}

\section{Introduction}
\label{sec:wave-direction-intro}

A single surface-following IMU contains directional information because the
horizontal and vertical components of wave-induced orbital motion are coupled.
The estimation problem nevertheless has two distinct ambiguities.  First,
horizontal motion identifies an \emph{axis}: the propagation plane is unchanged
when its representative vector is multiplied by $-1$.  Second, a physical
propagation direction requires choosing one of the two senses along that axis.
Treating those as one angle-estimation problem can hide a $180^\circ$ failure
behind otherwise accurate axial statistics.

The method studied here therefore separates the two tasks.  A frequency tracker
provides only a carrier for a quadrature Kalman estimator; both sine and cosine
coefficients are fitted, so the propagation axis does not depend on the
tracker's absolute phase.  The dominant eigenvector of the resulting horizontal
second moment gives an axial estimate and a linearity measure.  A separate
horizontal--vertical phase discriminator then resolves apparent travel sense.
The estimator is designed for real-time embedded use and reports uncertainty or
an unresolved state rather than forcing a direction when excitation is weak or
insufficiently linear.

The principal contributions are: (i) a phase-offset-invariant quadrature
formulation for the dominant wave axis; (ii) an explicit second-stage
travel-sense discriminator using the signed horizontal projection and vertical
orbital phase; (iii) confidence and gating rules that preserve the distinction
between an axis and a directed wavevector; and (iv) a simulation evaluation
that scores axial error and truth-referenced travel-sense correctness
separately.  The KalmANF carrier tracker used by the reported results is also
documented here; alternative compiled tracker policies are retained as
unevaluated appendices rather than being mixed into the direction claim.

\section{Signal Interfaces and Direction Conventions}
\label{sec:direction-interfaces}

The estimator consumes body-frame specific force, an attitude solution used to
remove roll--pitch gravity mixing, a projected bow heading, and a dominant
frequency estimate.  Its levelled horizontal frame is boat fixed: $+X$ is
forward, $+Y$ is starboard, and $+Z$ is up.  Angles formed with
$\operatorname{atan2}(Y,X)$ are therefore clockwise-positive when viewed from
above.  The method first reports a propagation-to vector; the conventional
waves-from direction is obtained by adding $180^\circ$.

\subsection{Attitude handoff}
\label{sec:proxy-handoff}
The direction stage does not estimate the vessel attitude itself.  In the
integrated implementation, a private gyro--accelerometer observer supplies
levelled tilt during startup and the navigation MEKF supplies attitude after
handoff.  The direction signals are resolved against the projected bow axis,
so arbitrary yaw of the startup tilt observer does not enter the horizontal
projection.  This interface lets the direction estimator remain a separable
measurement-side subsystem rather than a state of the navigation filter.

\input{w3d-frequency-tracking.tex-part}
\input{w3d-direction-methods.tex-part}
\input{w3d-wave-direction-config.tex-part}
\input{w3d-wave-direction-results.tex-part}

\section{Conclusion}
The separated estimator recovers a dominant wave axis from horizontal
quadratures and resolves its $180^\circ$ ambiguity with horizontal--vertical
orbital phase.  The committed simulation evidence shows small scenario-level
axial errors and a high truth-referenced travel-sense decision rate across the
eight tested directional seas.  The result should be read at that scope: it is
a dominant encounter-direction estimator, not a reconstruction of the full
directional spectrum, and it has not yet been physically referenced against a
directional wave instrument.  The most important next tests are crossing seas,
heading and vessel-speed sweeps, current-aware encounter correction, tracker
ablations, and external directional validation.

\appendices
\input{w3d-tracker-alternatives.tex-part}

\section*{Acknowledgments}
This work was prepared with the supervised assistance of artificial intelligence
tools, including ChatGPT, Claude and DeepSeek, which were employed to aid in
code generation, mathematical drafting, and editorial preparation. All
conceptual development, validation of results, and final revisions were
conducted independently by the author.

\section*{Disclosure Statement}
The author declares no conflicts of interest and received no funding for this
research.

\bibliographystyle{IEEEtran}
\bibliography{w3d,w3d-baselines}
\end{document}
